\section{家户、判牍与书写：地方社会的可见和不可见}

\subsection{户籍里的家庭并不完整}

户籍将人编为可征、可役、可管的单位，是金廷进入地方的必要工具。它也必然遗漏许多事情：寄居者何时离开，妇女怎样处置临时的债务，儿童何时加入劳动，雇工与佃户如何同主家协商，亲属是否分居而仍共享资源。史书中的“户口”让人口看上去可以被数清，却不是生活关系的原样复制。

战争使这种差异更大。某个成年男子被征发、被俘或随官署迁移，簿册可能仍以原来的家户称他；留在家中的人却要安排田土、收成、婚姻和老幼。金代墓志偶尔保留迁居、婚配与继承的线索，但它们写给有能力立石的家族。不能以精英的伦理语言推断所有家庭，也不能因贫者无名便说他们没有策略。材料要求我们把“国家登记了什么”与“家庭可能怎样维持”分开叙述。

\subsection{田界、债务与邻里}

土地的意义超过税基。田地是收获、居住、祖坟、水源、抵押和亲属关系相交的地方。战争或迁徙后，旧边界可能无人记得，新的占有者又需要某种认可；灌溉、道路和牲畜越过地界，也会使邻里关系进入争执。官府若受理案件，必须把复杂的使用关系翻译成文书类别；若当事人无法进入衙门，可能求助于亲属、里正、寺院或地方有声望者。我们缺少连续判牍，不能编造一套整齐的基层司法，却能看到为什么土地和债务总是国家分类的边缘。

借贷尤其不宜只写成剥削或救济。荒年借粮可能使一家渡过时日，也可能在收成不好时转成失地；商人借钱能够扩大交换，也可能在道路断绝后无力偿还。借约的利率、期限和担保取决于地域、当事人和保存下来的文本。没有金代全国性的契约档案，任何总论都必须克制。史家能做的是让一笔债务回到收获、运输、婚姻和战乱的具体风险中。

\subsection{法令为何要重复}

法令与条格保存的是国家说自己要如何治理。它们列出禁止、处罚、赏赐、程序和身份，因而能帮助我们理解行政词汇；它们不能自动显示一条命令已被人人理解或执行。重复颁令常意味着地方条件、官吏利益或信息传递使前一轮执行不完整。若只把重申当作严明，便会高估制度的穿透力；若只把它当作无效，也会忽略国家仍在试图重塑行为。

金代法制研究须同时面对文本缺失和后出汇编。一个保留下来的条目可能来自某次危机、某一地区或某项官僚整理，不能被当成全朝日常。宋方记录的金朝法令也经由对手的观察和政治需要筛选。用这些材料讨论司法时，应先问它规定什么、针对谁、由哪个机关发布；再问是否有地方文书、碑刻或其他叙事能显示其落地。没有第二层材料时，结论必须停在规范而非实践。

\subsection{女性的名字何时出现}

妇女在金代材料中常通过婚姻、母亲身份、节义叙述或功德题记出现。这样的出现并不等于她们只在这些关系中行动，而是保存机制倾向于这样记录她们。嫁资、寡居、收养与继承会影响财产和家户延续；战争与迁徙更可能使这些安排成为迫切问题。可惜，能够把一位女性的经济处境、亲属网络和个人选择同时保存下来的材料很少。

因此，本章不将后世熟悉的“贞节”叙事当作社会统计，也不以一两位有名女性证明普遍自由。比较可靠的说法是：家户在国家登记、财产关系和宗教纪念中被不断重新组织，女性承担的劳动、照护与财产处置常在文本边缘可见。边缘不等于无关；恰好因为这些工作连接了代际、土地和迁居，它们是理解战争如何进入日常的关键。

\subsection{学校、文集与仕途}

金廷经营学校、选举和礼制，是为了培养可以书写、审案和服从中枢秩序的人。地方士人进入这样的路径，可能获得仕进机会，也会把家族、地域和师承带进官场。文集和碑志保存较多的是成功者的声音，使我们容易把他们的价值判断误作全社会的共识。学校的存在不说明人人受教，科举名额不说明各地机会相同，文人的批评也不自动代表没有文书能力者的看法。

知识传播有具体成本。纸张、墨、书写者、藏书处和相对可通行的道路，决定一部书、一块碑或一封文书能否留下。金末的战争破坏这些条件，造成今天材料分布的不均。我们不能把这种不均写成文化不发达的证明，更不能以少数保存完好的都城资料覆盖边区。文化史应当同财政和交通史相连：谁提供纸与经费，谁保管档案，谁在战乱中能带走书籍，谁的记忆只能留在口耳之间。

\subsection{寺院的裁量空间}

寺院能提供礼拜、纪念和教育，也可能成为捐施、寄居和地方调解的节点。国家登记僧道、限制或保护寺产，显示朝廷将寺院纳入治理；实际上一座寺院与村落、商人、家族和官府的关系仍取决于地点与资源。功德碑所见的捐施者，说明某次工程已有可动用的网络；它没有告诉我们所有信众，也不能使我们把寺院想成脱离财政的避风港。

宗教文本的立场也需要辨认。僧传、碑记和灵验叙事重视传承、功德与神异，能照亮某些组织方式，却不等于现场记录。将它们与题记、建筑遗址、官府条目相互参照，才能知道一项叙述是在宣扬何种权威。材料的差异不是麻烦的附录，而是防止我们把地方社会写成只有国家一张脸的办法。
